\section{Experiments}

\subsection{Datasets and Fixed Protocol}

We evaluate on the 200-video SHOT test split \cite{zhu2023autoshot}, all 11 BBC
Planet Earth videos, and the 500-video ClipShots test split
\cite{tang2018clipshots}. SHOT is the primary short-video benchmark, BBC tests
long documentary footage, and ClipShots tests open-world transfer.

All headline numbers use the same controlled A1 logits and the B4 operating
point frozen by the validation-only procedure in Section III. Precision,
recall, and F1 are micro-averaged from transition-level TP/FP/FN counts.
Confidence intervals resample paired videos rather than individual frames.
They quantify test-set sampling uncertainty, not training-seed uncertainty.

\begin{table}[!t]
\caption{Fixed Reproducible Evaluation Protocol}
\label{tab:protocol}
\centering
\small
\renewcommand{\arraystretch}{1.08}
\begin{tabular}{L{1.45in}L{1.55in}}
\toprule
\textbf{Item} & \textbf{Setting} \\
\midrule
Training objective & BCE, one-hot labels \\
Training / data seed & 42 / 42 \\
Training budget & 20 epochs, batch size 512 \\
Sampling & At most 160 frames/video, 3 negatives/positive \\
Validation size & 200 videos \\
Validation selection & 5 source-stratified folds \\
Selection criterion & Mean macro-F1 across sources/folds \\
Selected temperature & \PaperTemperature \\
Gaussian smoothing & $\sigma=\PaperGaussianSigma$ \\
Decision threshold & $\theta=\PaperDeployThreshold$ \\
Matching tolerance & $\pm\PaperMatchTolerance$ frames \\
Bootstrap & 10,000 video resamples, seed 42 \\
\bottomrule
\end{tabular}
\end{table}

\subsection{Protocol-Matched Results}

Table~\ref{tab:main_results} compares configurations that share the controlled
Phase 2 protocol. B4 is selected exclusively by the frozen five-fold
validation manifest. Relative to A1, B4 improves SHOT by
\PaperASVTwoVsAOneShotPP{} F1 points and ClipShots by
\PaperASVTwoVsAOneClipPP{} points, while BBC is unchanged after rounding.
Temperature scaling alone preserves the validation-selected operating point;
Gaussian smoothing is the component that changes boundary decisions.

\begin{table}[!t]
\caption{Protocol-Matched F1 Scores (B4: Selected Method)}
\label{tab:main_results}
\centering
\small
\renewcommand{\arraystretch}{1.08}
\resizebox{\columnwidth}{!}{%
\begin{tabular}{L{2.15in}C{0.55in}C{0.55in}C{0.7in}}
\toprule
\textbf{Configuration} & \textbf{SHOT} & \textbf{BBC} & \textbf{ClipShots} \\
\midrule
\PaperProtocolMatchedRows
\bottomrule
\end{tabular}
}
\end{table}

Table~\ref{tab:confidence} reports the corresponding B4 operating points and
video-level bootstrap intervals. BBC has only 11 sampling units, so its
interval should be interpreted cautiously. ClipShots has high recall but lower
precision, indicating cross-domain false positives rather than a general
failure to detect transitions.

\begin{table}[!t]
\caption{B4 Operating Point and Video-Level Bootstrap 95\% Confidence Intervals}
\label{tab:confidence}
\centering
\scriptsize
\renewcommand{\arraystretch}{1.08}
\setlength{\tabcolsep}{2pt}
\begin{tabular}{L{0.48in}C{0.34in}C{0.95in}C{0.46in}C{0.46in}C{0.35in}}
\toprule
\textbf{Dataset} & \textbf{F1} & \textbf{F1 95\% CI} & \textbf{Precision} & \textbf{Recall} & \textbf{Videos} \\
\midrule
\PaperConfidenceRows
\bottomrule
\end{tabular}
\end{table}

Table~\ref{tab:paired_delta} directly bootstraps the per-video difference
between A1 at its validation-selected threshold and B4 at its frozen operating
point. The intervals support a positive B4 delta on SHOT and ClipShots. The BBC
interval contains zero, so we make no improvement claim for that dataset.

\begin{table}[!t]
\caption{Paired Video Bootstrap Difference, B4 Minus A1}
\label{tab:paired_delta}
\centering
\small
\renewcommand{\arraystretch}{1.08}
\resizebox{\columnwidth}{!}{%
\begin{tabular}{L{0.8in}C{0.6in}C{1.25in}C{0.65in}}
\toprule
\textbf{Dataset} & \textbf{$\Delta$F1} & \textbf{95\% CI} & \textbf{Excl. 0} \\
\midrule
\PaperPairedDeltaRows
\bottomrule
\end{tabular}
}
\end{table}

For context only, Table~\ref{tab:reported} lists numbers copied from the
original AutoShot and TransNetV2 reports. These values use their authors'
protocols and are not mixed into the controlled selection above.

\begin{table}[!t]
\caption{Literature-Reported F1 Scores (Context Only)}
\label{tab:reported}
\centering
\small
\renewcommand{\arraystretch}{1.08}
\resizebox{\columnwidth}{!}{%
\begin{tabular}{L{1.65in}C{0.55in}C{0.55in}C{0.7in}}
\toprule
\textbf{Method} & \textbf{SHOT} & \textbf{BBC} & \textbf{ClipShots} \\
\midrule
\PaperReportedBaselineRows
\bottomrule
\end{tabular}
}
\end{table}

\subsection{Calibration Diagnostics}

We evaluate frame probabilities on the combined 200-video validation split
using negative log-likelihood (NLL), Brier score, expected calibration error
with 15 equal-width bins (ECE), equal-mass adaptive ECE (AECE), and
class-balanced ECE (B-ECE). Table~\ref{tab:calibration_metrics} shows that
temperature scaling improves all diagnostics. The much larger B-ECE reveals
that aggregate frame metrics are dominated by negative frames. Validation
values remain diagnostic because the same split fits temperature.

\begin{table}[!t]
\caption{Frame-Level Validation Calibration Diagnostics}
\label{tab:calibration_metrics}
\centering
\scriptsize
\renewcommand{\arraystretch}{1.08}
\setlength{\tabcolsep}{2pt}
\begin{tabular}{L{0.85in}C{0.45in}C{0.45in}C{0.45in}C{0.45in}C{0.45in}}
\toprule
\textbf{Variant} & \textbf{NLL} & \textbf{Brier} & \textbf{ECE} & \textbf{AECE} & \textbf{B-ECE} \\
\midrule
\PaperCalibrationMetricRows
\bottomrule
\end{tabular}
\end{table}

\begin{figure}[!t]
\centering
\includegraphics[width=\columnwidth]{reliability_diagram.png}
\caption{Reliability diagram on the combined validation split. Temperature
scaling reduces aggregate error. Equal-mass bins expose sparsely populated
confidence regions.}
\label{fig:reliability}
\end{figure}

Table~\ref{tab:test_calibration} reports untouched-test calibration using the
temperature selected on validation. ClipShots has the largest B-ECE despite a
small aggregate ECE, consistent with domain-shifted false positives.

\begin{table}[!t]
\caption{Test Calibration at the Validation-Selected Temperature}
\label{tab:test_calibration}
\centering
\scriptsize
\renewcommand{\arraystretch}{1.08}
\setlength{\tabcolsep}{2pt}
\begin{tabular}{L{0.45in}C{0.47in}C{0.47in}C{0.47in}C{0.47in}C{0.47in}}
\toprule
\textbf{Dataset} & \textbf{NLL} & \textbf{Brier} & \textbf{ECE} & \textbf{AECE} & \textbf{B-ECE} \\
\midrule
\PaperTestCalibrationRows
\bottomrule
\end{tabular}
\end{table}

\subsection{Controlled Training Ablation}

Table~\ref{tab:ablation} expands the training-time comparison. A2, A3, and B1
show that Focal Loss and many-hot supervision do not improve the A1 control.
B5 combines these training changes with B4 and is numerically similar on SHOT
but slightly weaker on BBC and ClipShots. We therefore avoid attributing the
result to the more complex training objective. A0 is omitted because its
per-dataset best-threshold protocol is not a direct controlled comparison.

\begin{table}[!t]
\caption{Controlled Phase 2 Ablation (A0 Historical Baseline Excluded)}
\label{tab:ablation}
\centering
\scriptsize
\renewcommand{\arraystretch}{1.04}
\setlength{\tabcolsep}{2pt}
\begin{tabular}{L{1.7in}C{0.4in}C{0.4in}C{0.4in}}
\toprule
\textbf{Configuration} & \textbf{SHOT} & \textbf{BBC} & \textbf{ClipShots} \\
\midrule
\PaperControlledAblationRows
\bottomrule
\end{tabular}
\end{table}

\subsection{Efficiency and Reproducibility}

The backbone remains frozen, so only the Phase 2 head is optimized.
Table~\ref{tab:efficiency} reports parameter and run-artifact statistics.
The CPU post-processing measurement covers temperature scaling and Gaussian
smoothing on cached logits. End-to-end video decoding and backbone inference
are excluded. The original run artifact did not retain elapsed training time,
so we do not estimate or reconstruct it.

\begin{table}[!t]
\caption{Efficiency and Run-Artifact Statistics}
\label{tab:efficiency}
\centering
\small
\renewcommand{\arraystretch}{1.06}
\begin{tabular}{L{1.9in}L{1.05in}}
\toprule
\textbf{Item} & \textbf{Value} \\
\midrule
\PaperEfficiencyRows
\bottomrule
\end{tabular}
\end{table}

\subsection{ClipShots Transition-Type Recall}

We regenerate the ClipShots breakdown directly from the B4 logits and
\texttt{test.json}. One malformed annotation with reversed endpoints is
normalized by sorting the pair. A ground-truth interval is a cut when its
width is at most one frame and gradual otherwise. Predictions are
class-agnostic, so they are matched once and recall is grouped by the matched
ground-truth type. Type-specific precision and F1 are not reported because the
model does not predict a transition type.

\begin{table}[!t]
\caption{ClipShots Recall by Ground-Truth Transition Type}
\label{tab:clipshots_breakdown}
\centering
\small
\renewcommand{\arraystretch}{1.08}
\resizebox{\columnwidth}{!}{%
\begin{tabular}{L{0.8in}C{0.7in}C{0.65in}C{0.65in}C{0.65in}}
\toprule
\textbf{Type} & \textbf{GT} & \textbf{Matched} & \textbf{Missed} & \textbf{Recall} \\
\midrule
\PaperClipBreakdownRows
\bottomrule
\end{tabular}
}
\end{table}

The gradual-transition recall of 0.8089 is lower than the cut recall of
0.8991, but it is not the dominant source of the overall ClipShots gap.
The aggregate ClipShots precision is only 0.6498, indicating that open-world
effects and domain shift create many false-positive peaks.

\subsection{Training-Seed Robustness}
\label{sec:seed_robustness}

To quantify optimization variability, we retrained the BCE one-hot head
\PaperSeedCount{} times with training seeds 42, 43, and 44 while holding the
data seed, the training-video selection, and the sample cache identity fixed.
Each run selected its own post-processing configuration with the frozen
five-fold validation-only protocol of Section~\ref{sec:protocol} before any
test evaluation. Table~\ref{tab:seed_robustness} reports test means and
sample standard deviations. Seed-induced F1 variation is below 0.001 on SHOT
and BBC and below 0.005 on ClipShots, an order of magnitude smaller than the
video-sampling confidence intervals of Table~\ref{tab:confidence}; the
pipeline is therefore stable with respect to head-training randomness. The
absolute values differ from the headline B4 numbers because this study
retrains the head from scratch and re-selects its protocol per seed (seed 42
selected $\sigma{=}1.0$ and $\theta{=}0.30$), so the two tables answer
different questions: Table~\ref{tab:ablation} isolates component
contributions on the historical cache, whereas this study isolates
training-seed sensitivity end to end.

\begin{table}[!t]
\caption{Three-Seed Robustness of the End-to-End Pipeline (Mean $\pm$ Sample Standard Deviation Over Training Seeds 42--44)}
\label{tab:seed_robustness}
\centering
\small
\renewcommand{\arraystretch}{1.08}
\resizebox{\columnwidth}{!}{%
\begin{tabular}{L{0.85in}C{1.0in}C{1.0in}C{1.0in}}
\toprule
\textbf{Dataset} & \textbf{F1} & \textbf{Precision} & \textbf{Recall} \\
\midrule
\PaperSeedRobustnessRows
\bottomrule
\end{tabular}
}
\end{table}

\subsection{Archived Checkpoint Results and Limitations}

For auditability, Table~\ref{tab:archived} preserves the earlier deploy
checkpoint numbers. Only result JSON files remain for this checkpoint; its
test logits are absent. These numbers are therefore excluded from the headline
and confidence-interval analysis.

\begin{table}[!t]
\caption{Archived Deploy Checkpoint (Result JSON Only)}
\label{tab:archived}
\centering
\small
\renewcommand{\arraystretch}{1.08}
\resizebox{\columnwidth}{!}{%
\begin{tabular}{L{2.05in}C{0.55in}C{0.55in}C{0.7in}}
\toprule
\textbf{Archived protocol} & \textbf{SHOT} & \textbf{BBC} & \textbf{ClipShots} \\
\midrule
\PaperArchivedResultRows
\bottomrule
\end{tabular}
}
\end{table}

The study has four important limitations. First, the headline B4 head itself
was trained with one seed; the three-seed study of
Section~\ref{sec:seed_robustness} bounds optimization variability for the
retrained pipeline (F1 sample standard deviation at or below 0.005 on every
dataset), but it re-selects its protocol per seed and therefore does not
reproduce the historical B4 checkpoint bit for bit. Second,
the historical sample cache recorded 400 videos and 33,174 frames but not the
exact selected IDs; strict manifests fix this for future runs but cannot
retroactively prove the old cache contents. Third, validation calibration
diagnostics reuse the temperature-selection data. Fourth, the frozen
representation does not fully control cross-domain false positives on
ClipShots.
